\documentclass [varwidth]{standalone}
%scale[1.5]
%lmargin[20]
%vmargins[20]
\usepackage {texmaths-main}
\begin {document}
\setCompleteVersion {0}\par \begingroup \noindent \linewidth 340pt \begin {tblr}{ colspec={X[c,m,3]X[c,m,1.3]X[c,m,3]X[c,m,3]}, hlines,vlines, row{1}={font=\bfseries ,bg=\maintblrbgcolor }, rows={font=\small } } \toComplete {Objets}&\toComplete {Notation}&\toComplete {Représentation}&\toComplete {Remarques}\\ \toComplete {Le point A}&\toComplete {A}& \toComplete {\begin {tikzpicture}[baseline=(A.base)] \coordinate [label=above left:A](A)at(0,0); \draw node at(A){+}; \end {tikzpicture} } &\toComplete {Un point n'a pas d'épaisseur}\\ \toComplete {Le segment d'extrémités les points A et B}&\toComplete {[AB]}& \toComplete {\begin {tikzpicture}[baseline=(C.base)] \coordinate [label=above left:A](A)at(0,0); \coordinate (C)at(0,1); \coordinate [label=above:B](B)at(2,1); \tkzDrawLine [add=0 and 0](A,B) \draw node at(A){+}; \draw node at(B){+}; \end {tikzpicture} } &\toComplete {Un segment est fini}\\ \toComplete {La droite passant par les points A et B}&\toComplete {(AB)}& \toComplete {\begin {tikzpicture}[scale=0.8, baseline=(B.base)] \coordinate [label=above left:A](A)at(0,0); \coordinate [label=above:B](B)at(2,1); \tkzDrawLine [add=0.25 and 0.25](A,B) \draw node at(A){+}; \draw node at(B){+}; \end {tikzpicture} } &\toComplete {Une droite est infinie}\\ \toComplete {La demi-droite d'origine A et passant par B}&\toComplete {[AB)}& \toComplete {\begin {tikzpicture}[scale=0.8, baseline=(B.base)] \coordinate [label=above left:A](A)at(0,0); \coordinate [label=above:B](B)at(2,1); \tkzDrawLine [add=0 and 0.35](A,B) \draw node at(A){+}; \draw node at(B){+}; \end {tikzpicture} } &\toComplete {Une demi-droite est infinie d'un côté} \end {tblr}\par \endgroup 
\end {document}
